\section{Motivation}

The increasing importance of open science and reproducibility has 
led to a growing need for effective management and reuse of digital 
research outputs, including data, software, and computational 
workflows. The FAIR principles — Findable, Accessible, Interoperable, 
and Reusable — were introduced by Wilkinson et al. \cite{wilkinson2016fair} 
to address these challenges by providing guidelines that enhance the 
discoverability and reuse of research artifacts. A central goal of 
the FAIR principles is machine-actionability — the ability of 
automated agents to discover, evaluate, and act on digital objects 
without prior knowledge of them.

While initially developed for research data, the FAIR principles are 
equally applicable to research software, which plays a central role 
in modern scientific processes. The FAIR for Research Software 
(FAIR4RS) principles \cite{barker2022fair4rs} extend and adapt the 
original framework to address the specific characteristics of 
software, such as executability, versioning, and dependency 
management. Ensuring that software is FAIR is essential for enabling 
reproducibility, transparency, and collaboration.

However, in practice, many research software repositories lack 
essential elements including structured metadata, clear licensing 
information, persistent identifiers, and comprehensive documentation. 
Assessing the FAIRness of research software at scale remains a 
challenging task, particularly for indicators that depend on the 
quality and completeness of software documentation 
\cite{chuehong2025frsm}.

Existing FAIRness assessment approaches primarily rely on rule-based 
checks that evaluate the presence of metadata files, licenses, 
persistent identifiers, and other objectively measurable indicators 
\cite{fuji2021}. While these approaches provide valuable automated 
assessments, they often struggle to evaluate qualitative aspects such 
as installation guidance, usage instructions, and documentation 
quality — factors that directly influence software reuse and 
reproducibility \cite{chuehong2025frsm}.

Recent advances in Large Language Models (LLMs) offer new 
opportunities for analyzing natural language documentation and 
supporting FAIRness assessment beyond traditional rule-based methods. 
By interpreting README files and other documentation artifacts, LLMs have the potential to evaluate aspects that are difficult to capture using deterministic rules alone.

This thesis investigates how LLM-based documentation assessment can
complement rule-based FAIRness evaluation for research software
through a hybrid assessment approach that combines automated
file-based checks with LLM-supported analysis of repository
documentation. The goal is to investigate whether LLM-based assessment
can capture qualitative FAIR-related documentation characteristics
that are not addressed by rule-based methods alone.

\section{Problem Statement}

Current automated methods for evaluating the FAIRness of research
software repositories are well suited to assessing explicit and
machine-detectable properties, but are more limited when information
is expressed in unstructured documentation. Existing automated
FAIRness assessment approaches commonly rely on deterministic checks
that evaluate the presence of metadata files, licenses, persistent
identifiers, and other objectively measurable indicators
\cite{fuji2021}. However, such approaches are not designed to assess
qualitative characteristics of software documentation, such as the
clarity and completeness of installation guidance, usage
instructions, or dependency information.

As noted by Chue Hong et al. \cite{chuehong2025frsm}, automated tools
can verify the existence of metadata, but presence alone does not
establish its quality, relevance, or usefulness.

Research on Mining Software Repositories has shown that README files 
are the primary documentation artifact for GitHub repositories and 
contain rich information about a project's purpose, usage, and 
quality \cite{prana2019readme}. However, the systematic integration of documentation-based information into automated FAIRness assessment workflows remains limited. 
Rule-based checks can readily verify whether a README file exists, but
presence-based checks alone do not determine whether it adequately
describes installation procedures, usage examples, or software
dependencies. Treude et al. 
\cite{treude2020beyond} further demonstrate that documentation 
quality goes beyond mere presence — dimensions such as clarity, 
understandability, and structure require qualitative assessment that 
rule-based approaches cannot provide.

The importance of documentation quality extends beyond academic 
frameworks. Imani et al. \cite{imani2024documentation} found that 
over 91\% of industry practitioners rely on OSS documentation for 
software adoption decisions, and that the top reported challenge is 
not the absence of documentation but the difficulty of finding 
relevant information within it. This practitioner evidence reinforces 
the need for assessment methods that evaluate not just the presence 
but the quality and completeness of documentation content.

Consequently, there is a need to investigate whether a hybrid approach
that combines rule-based structural checks with LLM-based qualitative
assessment can extend automated FAIRness evaluation to documentation
characteristics that are not captured by rule-based methods alone.

\section{Research Objective and Research Questions}
The primary objective of this thesis is to design, implement, and
evaluate a prototype system for the automated assessment of FAIRness
in research software repositories hosted on GitHub. The system
combines rule-based indicators derived from the FAIR4RS principles
\cite{barker2022fair4rs} with an LLM-based documentation assessment
component. The evaluation investigates both the extent to which
LLM-based assessment complements deterministic structural checks by
capturing qualitative documentation characteristics and the
reliability and failure modes of such assessment across two evaluated
LLM configurations.

The following research questions guide the investigation:

\begin{description}
    \item[RQ1] To what extent can LLM-based assessment complement
rule-based FAIRness assessment by capturing qualitative documentation
characteristics of research software repositories?
    
   \item[RQ2] How reliably can LLM-based methods assess qualitative
documentation dimensions of research software repositories, and what
failure modes affect their assessments?

\end{description}

The system is implemented as a modular Python prototype focused on 
GitHub-hosted repositories. Extensibility to other platforms such as 
Zenodo or Codeberg is considered out of scope for this thesis but is 
noted as a direction for future work.

\section{Thesis Structure}

The remainder of this thesis is organized as follows.

Chapter 2 reviews the relevant literature, covering the FAIR 
principles and their adaptation for research software, existing 
automated FAIRness assessment tools, research on software 
documentation quality, and recent advances in LLM-based document 
analysis.

Chapter 3 describes the proposed hybrid assessment approach, 
including the dataset collection methodology, the rule-based
component, the LLM-based documentation assessment component, and
the scoring and output design.

Chapter 4 describes the implementation of the proposed framework,
including the repository processing pipeline, rule-based assessment
component, LLM integration, generation of hybrid FAIRness profiles,
and the web-based prototype for exploring saved results and assessing
public GitHub repositories.

Chapter 5 presents the evaluation strategy and results, including
ground-truth validation, quantitative comparison of the evaluated
LLM configurations against manual annotations, analysis of
dimension-level scoring behaviour and failure patterns, and
comparison with existing FAIRness assessment tools.

Chapter 6 discusses the findings in relation to the research
questions, considers the implications and limitations of the
approach, and outlines directions for future work.

Chapter 7 concludes the thesis by summarising the main contributions
and findings.